\SourceLeaf{494}%
sind die Theodolithen, wie sie der englische Künstler \textsc{Troughton} verfertigt, am bequemsten und brauchbarsten. Der Prof. \textsc{Schumacher} hat die Gefälligkeit gehabt, mir den seinigen zu diesem Behuf für den Einkaufspreis wieder abzutreten, indem er bis zu der Zeit, wo er selbst wieder in diesem Jahre einen solchen nöthig hat, einen andern aus England zu erhalten hofft. — Manche Winkelmessungen werden sich nicht gut anders, als bei Nacht durch Argandsche Lampen mit parabolischem Reverberes, die auf sehr grosse Weiten sichtbar gemacht werden können, anstellen lassen. Damit es in vorkommenden Fällen daran nicht fehle, habe ich mehrere solche Reverberes bei einem sehr geschickten Künstler, der auch ähnliche Arbeiten für die preussischen Messungen geliefert hat, dem Hof-Mechanikus \textsc{Körner} in Jena, bestellt, der solche auf Ostern zu vollenden versprochen hat. Mehrerer anderer in Voraus besprochener Apparate hier nicht zu gedenken.

Nach allen diesen Vorkehrungen glaube ich nun die Arbeiten im bevorstehenden Frühjahr anfangen zu können. Wenn die Künstler alle ihre Versprechungen gehörig inne halten, wird das Werk dann immer rasch fortschreiten können. Auf alle Fälle aber werden die mannigfaltigen Vorarbeiten, Bereisungen der Gegenden, Aussuchung der Stationen, Erbauungen von Signalthürmen, vorläufige Messungen u. dergl. — (der vorhinerwähnte Theodolith von \textsc{Troughton} ist bereits in meinen Händen) — erst mehrere Monate ausfüllen, und selbst den widrigen Fall angenommen, dass, nachdem die Vorarbeiten so weit vorgeschritten wären, dass die Messungen selbst beginnen könnten, doch der erwartete grössere \textsc{Reichenbach}'sche Theodolith noch nicht abgeliefert wäre, würde die Arbeit nicht still zu stehen brauchen, da ich im Nothfall mit denjenigen Winkelmessungen anfangen könnte, wozu die anderweitigen Hülfsmittel der Sternwarte zureichen. Ja im allerschlimmsten Fall, der hoffentlich nicht eintreten wird, dass die Ablieferung des Theodolithen noch länger verzögert würde, würde ich nur in der Ordnung der Arbeiten die Abänderung eintreten lassen, dass ich den Rest des Sommers zu den astronomischen Arbeiten mit dem Zenithsector hier in Göttingen verwende, die sonst, wenn der ganze Sommer dem Triangulationsgeschäft gewidmet werden kann, einer spätern Zeit vorbehalten bleiben.

Göttingen, den 20. März 1821.
\begin{flushright}C. F. \textsc{Gauss}.\end{flushright}

\par\bigskip\Needspace{13\baselineskip}
\pdfbookmark[0]{Druckfehler.}{druckfehler}
\begin{center}\large DRUCKFEHLER.\end{center}
\par\noindent Seite 296 statt:
\begin{center}
\begin{tabular}{c|c|c|c}
\(Q+q\)&\(P+p\)&\(\log m\)&\(k\)\\
&&\(+\)&\\
\(52^{0}\ 30'\)&\(52^{0}\ 32'\ 1''78428\)&&\(0''006\)
\end{tabular}
\end{center}
\noindent lies:
\begin{center}
\begin{tabular}{c|c|c|c}
\(Q+q\)&\(P+p\)&\(\log m\)&\(k\)\\
&&\(-\)&\\
\(52^{0}\ 30'\)&\(52^{0}\ 32'\ 1''78428\)&&\(0''006\)
\end{tabular}
\end{center}
